\fig{fig:w2:energyflux}{左侧弦通过截面向右侧传力；右行波的横向力与本地点速度同向}{\begin{tikzpicture}

\draw[world,domain=-2:2,samples=80]plot(\x,{-.45*\x+.08*\x*\x});
\fill(0,0)circle(.06)node[below right]{$C$};\draw[guide](0,-1.4)--(0,1.7);\node[above]at(0,1.7){截面$x$};
\draw[force](0,0)--(-1.5,.675)node[above left]{$\boldsymbol F_{L\to R}$};\draw[flow](.3,0)--(.3,1.05)node[right]{$v_y=-u y_x>0$};
\draw[flow](1,-1.35)--(2.5,-1.35)node[right]{$P>0$};\node at(-1.8,-1.2){左侧};\node at(2.7,-.8){右侧};

\end{tikzpicture}}
图\ref{fig:w2:energyflux}取$y_x<0$的一个瞬间。左侧对右侧的张力指向左上，其竖直分量为$-\mathcal T y_x>0$；右侧端点此刻也向上，所以左侧向右侧做正功。若$y_x>0$，力与速度的竖直分量同时变负，功率仍为正。这里计算通过截面的能流，故应取左侧作用于右侧的力，不能误取反作用力。

由弦元的牛顿方程，还可直接检验局部能量守恒。令$u_E=\mu_\ell y_t^2/2+\mathcal T y_x^2/2$，则
\begin{align*}
\partial_tu_E
&=\frac{\mu_\ell}{2}(2y_ty_{tt})+\frac{\mathcal T}{2}(2y_xy_{xt})\\
&=\mathcal T(y_ty_{xx}+y_xy_{xt})
=\mathcal T\partial_x(y_ty_x)=-\partial_xP.
\end{align*}
第二行先代入$\mu_\ell y_{tt}=\mathcal T y_{xx}$，再用$\partial_x(y_ty_x)=y_{tx}y_x+y_ty_{xx}$与混合偏导$y_{tx}=y_{xt}$合并为乘积导数。
因此一小段弦内能量的增加等于左端流入功率减去右端流出功率。对单一右行波还有$P=u u_E$；对含入射与反射的叠加场，能量密度依然非负，净功率却可能为零，不能把总能量密度乘$u$一律视为净能流。
